\section{Findings and Discussion}\label{sec:findings}

\subsection{Key Findings}

The central finding is that \textbf{the conversational system does what the oracle asks, while
the baseline only looks tidier}. The conversational system wins operation appropriateness with
disjoint intervals ($3.94\pm0.31$ vs.\ $2.73\pm0.35$) and trends higher on intent alignment and
oracle rating; the baseline's higher coherence and label scores are an artefact of producing
almost no structure (Finding~2). The remaining findings are tagged by evidential strength:
disjoint intervals (CI-clean), a direction with overlapping intervals (trend), a proportion or
distribution, a counted defect, or a qualitative observation.

\begin{description}[leftmargin=1.4em,style=nextline]

  \item[F1: The explanation agent explains the \emph{wrong film} about half the time
    \textnormal{(counted defect, $n=89$).}] Across the 89 ``why is film $X$ here'' turns,
    43 ($48\%$) answered about a different title: a question about \emph{Don't Look Up} was
    answered with ``Playdate'', and one about \emph{Mean Girls} with ``Friends with Benefits''.
    This systematic grounding failure is the hard number behind the lowest judge score in
    the study, explanation quality $=1.79\pm0.31$.

  \item[F2: Two opposite clustering failure modes \textnormal{(distributional).}] The baseline
    ends with two or fewer clusters in $87\%$ of sessions (26/30; maximum 7), a chronic
    under-segmentation. The conversational system usually produces a handful of clusters but
    occasionally \emph{explodes}, reaching 17--18 clusters in four sessions (maximum 18). The
    baseline's coherence and label-accuracy edge is therefore a coarseness artefact, and cluster
    counts must be read as distributions, not means.

  \item[F3: A capability cliff between structured and abstract intents \textnormal{(trend,
    small $n$).}] Per-persona oracle ratings run from 2.0 for abstract goals (a satirical
    \emph{tone}, the \emph{New Hollywood} movement, the \emph{heist} sub-genre) up to 4.67 for
    metadata-grounded ones (a genre fan, a rater). The system is a better \emph{clusterer} of
    well-scoped attributes than a \emph{grasper} of fuzzy semantic concepts.

  \item[F4: On a genuinely abstract intent, the baseline \emph{wins} \textnormal{(trend, $n=3$
    each).}] For the \texttt{newhollywood} persona the conversational system scored 2.0 against
    the baseline's 5.0. A single holistic pass can sometimes capture a fuzzy film-movement concept
    that the operation-constrained pipeline, forced to express everything as discrete navigation
    steps, cannot. This complicates the ``conversational is better'' story honestly.

  \item[F5: The conversational system completes more and misbehaves less
    \textnormal{(proportion).}] It finishes the intended trajectory in $49\%$ of sessions and is
    stopped for misbehaviour in $34\%$; the baseline finishes in $37\%$ and misbehaves in $53\%$,
    a majority of its sessions.

  \item[F6: Going multi-agent costs essentially the same per turn as one model call
    \textnormal{(CI-tight).}] Cost per turn is \$0.0126 vs.\ \$0.0123 (per session \$0.169 vs.\
    \$0.153). The agentic overhead is near-zero because about a quarter of conversational turns
    are cheap non-operational moves; operations per turn are $0.74$ vs.\ $0.93$. The usual cost
    objection to ``going agentic'' does not hold here.

  \item[F7: Oracle satisfaction is bimodal \textnormal{(distributional).}] Ratings cluster toward
    the extremes (1--2 and 4--5) with very few middle values, so a session tends to be a clear hit
    or a clear miss and the mean rating is a poor summary; the distribution is the informative
    object.

  \item[F8: Oracle ratings are lenient and only weakly tied to cluster quality
    \textnormal{(qualitative + distributional).}] The top-rated conversational session (oracle
    $=5$) had a judged clustering coherence of 2, visibly looped, and returned two groups when
    three were requested (Section~\ref{sec:results}). The oracle rewards responsiveness and
    recovery more than the quality of the final partition, a validity caveat, and itself an
    informative disagreement between oracle and judge.

  \item[F9: Cluster labels concatenate and bloat with drill-down depth
    \textnormal{(qualitative).}] The labeller \emph{prepends} to inherited labels instead of
    renaming, yielding names such as \textsf{Light-hearted Romantic Comedies Asian Romantic
    Comedies Family Fantasy Comedies}. Depth compounds the problem rather than refining it.

  \item[F10: The requested number of clusters is frequently not honoured
    \textnormal{(qualitative).}] ``Split into 3 groups'' returned two groups in one session and
    four in another. The split count the oracle asks for is treated as a hint rather than a
    constraint.

\end{description}

Read together, these findings tell a coherent story. The conversational architecture's value is
\emph{control fidelity}: choosing the right operation and following the oracle's trajectory, not
intrinsic cluster quality, where it is no better and sometimes worse than a single model (F2, F4).
Its sharpest weaknesses (F1, F9, F10) are localised, mechanical defects in individual agents
rather than failures of the overall loop, which makes them tractable to fix. And the evaluation
surfaces a measurement caveat of its own: the oracle is a lenient judge of structure (F7, F8), so
oracle rating and judge dimensions should be read together rather than either alone.

\subsection{Limitations}

The strongest constraint is the \textbf{small, forced number of experiments}. With roughly 30
sessions per condition only one comparison, operation appropriateness, reaches disjoint confidence
intervals; every other difference is a trend with overlapping intervals and should be treated as
suggestive, not established. The per-persona findings (F3, F4) rest on as few as three sessions,
where the confidence interval is necessarily wide. Only \emph{two} conditions were run, on a
single compact 288-film catalogue, so neither catalogue scale nor a fuller ablation of the
architecture has been probed. Finally, both the user and the grader are language models: these
results characterise agent-to-agent dynamics, and the leniency of the simulated oracle (F8) means
oracle ratings in particular should not be read as human-user satisfaction.

\subsection{Future Work}

\paragraph{Close the local defects.} The sharpest weaknesses are mechanical and should be fixed
first: ground the explanation agent to the exact film it is asked about (F1), make the labeller
\emph{rename} rather than prepend (F9), and treat the requested split count as a constraint (F10).

\paragraph{Strengthen the intent layer.} The intent layer is today a bottleneck on multi-step
instructions and would benefit from more robust parsing of compound requests, possibly with a
verification pass over the parsed actions.

\paragraph{Evaluation.} More sessions and seeds would tighten the intervals enough to promote
today's trends to firm results, and a human-oracle study would test how far the agent-to-agent
findings transfer to real users.

\paragraph{New domains.} Because the core loop, in which the oracle proposes a grouping and the
system partitions and refines, is not specific to films, the system is a natural candidate for
domain-agnostic catalogues such as books, papers, or products, once the embedding pipeline and the
movie-specific prompts are made configurable. Exercising it on catalogues of varying size and
attribute density would surface the edge cases a single curated catalogue cannot.
